\section{The 30-Mile Standard: National Coverage Findings}

\subsection{Overall Coverage}

Of 3,105 continental US counties, 1,595 (51.4\%) have their population centroid within 30 miles of an interstate on-ramp. This means 1,510 counties (48.6\%) fall outside the 30-mile threshold --- a surprisingly high fraction given that the interstate system spans 48,800 miles.

However, the distribution is heavily population-weighted. Counties within the 30-mile threshold tend to be more densely populated (containing most major metropolitan areas), while gap counties are predominantly rural with low population density. Table~\ref{tab:coverage} presents the full coverage statistics.

\begin{table}[h!]
\centering
\caption{Interstate Coverage of Continental US Counties, 2023}
\label{tab:coverage}
\begin{tabular}{lrrrrr}
\toprule
\textbf{Threshold} & \textbf{Counties} & \textbf{\%} & \textbf{Population} & \textbf{\% of US} & \textbf{Land (sq mi)} \\
\midrule
Within 20 miles & 1,008 & 32.5\% & 204,130,890 & 61.9\% & 529,420 \\
Within 30 miles & 1,595 & 51.4\% & 262,331,636 & 79.6\% & 1,076,380 \\
Within 50 miles & 2,361 & 76.1\% & 295,866,227 & 89.8\% & 2,010,000 \\
\textbf{Beyond 30 miles} & \textbf{1,510} & \textbf{48.6\%} & \textbf{66,513,310} & \textbf{20.4\%} & \textbf{2,118,000} \\
\bottomrule
\end{tabular}
\end{table}

The I2.0 coverage target --- 99\% of the US population within 30 miles of an interstate on-ramp via the combined T1/T2/T3 network --- requires bringing an additional 19.3\% of the population (63.7 million people) within threshold. The 0.6\% remaining beyond threshold after the full I2.0 program corresponds to the most geographically isolated communities (deep desert, high mountain terrain) that are candidates for rural access spurs rather than full interstate corridors.

\subsection{Geographic Distribution of Gap Counties}

Figure~\ref{fig:gap-map} shows the geographic distribution of the 1,510 gap counties. The distribution is strongly non-uniform, with four visible clusters that correspond to the four gap zones analyzed in Section~\ref{sec:taxonomy}.

\begin{figure}[h!]
\centering
\includegraphics[width=\textwidth]{figures/gap-counties-map.pdf}
\caption{Continental US gap counties (centroid $> 30$ miles from nearest interstate on-ramp). Color intensity proportional to gap distance. Four gap zones visible: Northern Tier (upper left), Appalachians (center-east), Rural Gulf South (south-center), Rural West (right).}
\label{fig:gap-map}
\end{figure}

\subsection{State-Level Results}

Table~\ref{tab:states} shows the 15 states with the largest gap populations, which collectively account for 73\% of the total gap.

\begin{table}[h!]
\centering
\caption{Top 15 States by Gap Population (Counties $> 30$ Miles from Interstate)}
\label{tab:states}
\begin{tabular}{lrrr}
\toprule
\textbf{State} & \textbf{Gap Counties} & \textbf{Gap Population} & \textbf{\% of State Pop} \\
\midrule
California & 34 & 12,177,286 & 30\% \\
Arizona & 14 & 7,155,601 & 94\%$^\dagger$ \\
Louisiana & 45 & 2,631,301 & 57\% \\
Tennessee & 42 & 1,660,090 & 23\% \\
South Carolina & 21 & 1,844,185 & 34\% \\
Washington & 24 & 1,360,949 & 17\% \\
Utah & 20 & 1,353,166 & 38\% \\
Missouri & 52 & 1,339,306 & 21\% \\
Illinois & 44 & 1,442,540 & 11\% \\
Pennsylvania & 13 & 1,129,530 & 9\% \\
New Mexico & 21 & 874,061 & 40\% \\
Michigan & 43 & 1,388,641 & 13\% \\
Ohio & 17 & 835,741 & 7\% \\
North Carolina & 24 & 1,010,123 & 9\% \\
Kentucky & 66 & 1,698,632 & 37\% \\
\midrule
\textbf{Top 15 subtotal} & 480 & 40,860,152 & 20\% \\
\textbf{All gap counties} & 1,510 & 66,513,310 & 20.4\% \\
\bottomrule
\multicolumn{4}{l}{$^\dagger$ Arizona's high percentage reflects the county-centroid effect for large desert counties.} \\
\multicolumn{4}{l}{Tucson is on I-10; the county centroid falls in the eastern Sonoran Desert.} \\
\end{tabular}
\end{table}

California's high gap population (12.2 million) reflects the large geographic extent of inland counties (San Bernardino County is the largest county by area in the continental US at 20,000 sq mi; its centroid is in the Mojave Desert, far from I-10 or I-40). Louisiana's 45 gap counties reflect the relative sparsity of the interstate network in south Louisiana and the northern parishes.

\subsection{The Worst Individual Gaps}

The worst continental US gaps (by distance) are in northern Maine: Aroostook County (261 miles, 67,237 people), Washington County (202 miles, 31,096 people), and Piscataquis County (202 miles, 16,936 people). Maine has no interstate north of I-95 at Augusta, leaving the entire northern two-thirds of the state without interstate access. The Northern Tier cluster in North Dakota and Montana includes Divide County, ND (186 miles, 2,195 people) and Sheridan County, MT (183 miles, 3,700 people).

These worst-gap counties represent the geographic argument for the Northern Tier interstate: not just economic development, but healthcare access, agricultural logistics, and disaster evacuation for communities with no current alternative.

\subsection{The Economic Opportunity Correlation}

We score each county on the C3 (Economic Opportunity Access) dimension of the ROUTE scoring framework, using BEA CAINC4 2022 county GDP-per-capita data. Gap counties (centroid $> 30$ miles from interstate) have a mean C3 score of 5.8, compared to 4.1 for non-gap counties --- a 41\% difference, statistically significant at $p < 0.001$ (Mann-Whitney U-test).

A C3 score of 5.8 corresponds to buffer GDP per capita at approximately 77\% of the national average; a score of 4.1 corresponds to approximately 88\%. Gap counties are not just far from interstates --- they are systematically below the national economic average in a way that compounds their geographic isolation.

This finding has a direct policy implication: closing the coverage gap does not only provide transportation access. It connects below-average economic regions to the national economy in a way that the coverage measure alone does not capture. The investment case for priority corridors is therefore stronger than the pure coverage numbers suggest.
